%\section*{Calculations}

%\subsection*{Mass fraction of TC in the gel and xerogel}

%The mass fraction of TC in the gel is calculated as
%\begin{equation}
%    \text{mass fraction of TC in the gel}
%    = \frac{m_{\mathrm{chitosan,aliquot}}}{m_{\mathrm{gel,total}}}.
%    \label{eq:tc-gel-mass-fraction}
%\end{equation}
%
%The mass fraction of TC in the xerogel is calculated as
%\begin{equation}
%    \text{mass fraction of TC in the xerogel}
%    = \frac{m_{\mathrm{gel,taken}} \times \text{mass fraction of TC in the gel}}{m_{\mathrm{xerogel}}}.
%    \label{eq:tc-xerogel-mass-fraction}
%\end{equation}
%
%\subsection*{Thiol content per milligram of TC}
%
%For a gel sample, the thiol content per milligram of TC is calculated as
%\begin{equation}
%    \text{thiol content per mg TC in the gel}
%    = \frac{c_{\mathrm{thiol}}\,V_{\mathrm{assay}}\,D}
%           {m_{\mathrm{gel}} \times \text{mass fraction of TC in the gel}}.
%    \label{eq:thiol-content-gel}
%\end{equation}
%
%For a xerogel sample, the corresponding calculation is
%\begin{equation}
%    \text{thiol content per mg TC in the xerogel}
%    = \frac{c_{\mathrm{thiol}}\,V_{\mathrm{assay}}\,D}
%           {m_{\mathrm{xerogel}} \times \text{mass fraction of TC in the xerogel}}.
%    \label{eq:thiol-content-xerogel}
%\end{equation}
%
%Here, $c_{\mathrm{thiol}}$ is the thiol concentration determined using Ellman's assay,
%$V_{\mathrm{assay}}$ is the assay volume, $D$ is the dilution factor, and $m$ denotes the
%corresponding sample mass. The resulting value gives the amount of thiol per milligram of TC.


\section*{Plots}

\subsection*{Ellman's assay optimizations}

\begin{figure}[H]
  \centering
  \begin{subfigure}[b]{0.48\textwidth}
    \centering
    \includegraphics[width=\linewidth]{Images/plots/ellmans/Gel2.pdf}
    \caption{Comparison of the standard Ellman's assay protocol with adapted methods for thiolated chitosan reported in the literature.}
    \label{fig:asc1}
  \end{subfigure}
  \hfill
  \begin{subfigure}[b]{0.48\textwidth}
    \centering
    \includegraphics[width=\linewidth]{Images/plots/ellmans/Timepoints_richtig.pdf}
    \caption{Amount of thiol groups detected using Ellman's assay at different time points}
    \label{fig:asc2}
  \end{subfigure}
  \caption*{\textbf{Supplementary Figure S1.} Ellman's assay optimization tries.}
  \label{fig:main}
\end{figure}


 \begin{figure}[htbp]
    \centering
    \begin{subfigure}[c]{0.38\textwidth}
    \includegraphics[width=\linewidth]{Images/results/chit-asc-with-anno.png}
    \label{fig:chit-asc-pic}
    \caption{Concentration-dependent color change following the addition of ascorbic acid to the Ellman's assay.}
    \end{subfigure}
    \hfill
      \begin{subfigure}[c]{0.45\textwidth}
    \centering
    \includegraphics[width=\linewidth]{Images/plots/ellmans/Ansätze.pdf}
    \caption{Quantification of thiol-group content in samples with Ascorbic acid using the Ellman's assay.}
    \label{fig:alter1}
  \end{subfigure}
  \caption{\textbf{Supplementary Figure S2.} Addition of ascorbic acid to the Ellman's assay mixtures. The excess in the two samples on the left was calculated based on the detected thiol content in the xerogel 4 (e.g., 1.2 $\times$ 24~$\mu$mol/mg), whereas the excess in the sample on the right (literature) was calculated based on the expected thiol group content, reported by the literature (1.5 $\times$ 140~$\mu$mol/mg).}
  \end{figure}

\subsection*{Encapsulation efficiency}

\begin{figure}[htbp]
  \begin{subfigure}[c]{\textwidth}
    \centering
    \includegraphics[width=0.39\linewidth]{Images/plots/EE-MalvsOhneMal.pdf}
  \caption*{Encapsulation efficiency of liposome formulations with- vs without maleimide.}
  \label{fig:EE-malohneaml}
  \end{subfigure}
  \hfill
    \begin{subfigure}[c]{\textwidth}
      \centering
    \includegraphics[width=0.48\linewidth]{Images/plots/EE-all.pdf}
    \caption*{Encapsulation efficiency of all liposomes.}
\label{fig:EE-all}
  \end{subfigure}
  \label{fig:EE-sup}
  \caption*{\textbf{Supplementary Figure S3.} Encapsulation efficiency of liposome formulations with- vs without maleimide.}
\end{figure}


\begin{table}[ht]
  \centering
  \small
  \caption{Amicon$^\text{\textregistered}$ centrifugation table}
  \label{tab:Amicon}
  \setlength{\extrarowheight}{0.1cm}
  \begin{tabular}{c l c c c}
    \toprule
    Centrifugation step & Solution & RPM (g) & Duration (min) & Eluate (ml) \\
    \midrule
    1  & Chitosan          & 3986 & 15 & $\sim$1    \\
    2  & Chitosan          & 4800 & 15 & $\sim$1    \\
    3  & Chitosan          & 4800 & 15 & $\sim$1.5  \\
    4  & Chitosan-HCl NaCl & 4000 & 25 & $\sim$2    \\
    5  & Chitosan-HCl NaCl & 4000 & 40 & $\sim$2.5  \\
    6  & Chitosan-HCl NaCl & 4000 & 40 & 3          \\
    7  & Chitosan-HCl NaCl & 4000 & 40 & 3          \\
    8  & Chitosan-HCl NaCl & 4000 & 40 & 3.5        \\
    9  & Chitosan-HCl NaCl & 4000 & 40 & 3.75       \\
    10 & Chitosan-HCl NaCl & 4000 & 35 & 3.5        \\
    11 & Chitosan-HCl NaCl & 4000 & 30 & 3.9        \\
    12 & Chitosan-HCl NaCl & 4000 & 30 & 3.5        \\
    13 & Chitosan-HCl NaCl & 4000 & 30 & 3.5        \\
    14 & Chitosan-HCl NaCl & 4000 & 10 & 0.5        \\
    \bottomrule
  \end{tabular}
\end{table}